\documentclass[conference]{IEEEtran}
\IEEEoverridecommandlockouts
% The preceding line is only needed to identify funding in the first footnote. If that is unneeded, please comment it out.
\usepackage{cite}
\usepackage{amsmath,amssymb,amsfonts}
\usepackage{algorithm, algorithmicx}
\usepackage[noend]{algpseudocode}
\usepackage{graphicx}
\usepackage{textcomp}
\usepackage{xcolor}
% Add by author
\usepackage{cite}
\usepackage{amsmath}
% \usepackage{fontspec}


\def\BibTeX{{\rm B\kern-.05em{\sc i\kern-.025em b}\kern-.08em
    T\kern-.1667em\lower.7ex\hbox{E}\kern-.125emX}}
\begin{document}

\title{Visual Fidelity-Guaranteed Sampling for Large-scale Trajectory Data Visualization
}

\author{\IEEEauthorblockN{1\textsuperscript{st} Given Name Surname}
\IEEEauthorblockA{\textit{dept. name of organization (of Aff.)} \\
\textit{name of organization (of Aff.)}\\
City, Country \\
email address or ORCID}
\and
\IEEEauthorblockN{2\textsuperscript{nd} Given Name Surname}
\IEEEauthorblockA{\textit{dept. name of organization (of Aff.)} \\
\textit{name of organization (of Aff.)}\\
City, Country \\
email address or ORCID}

}

\maketitle

% Defined by authors--------------------------------------------
\newtheorem{problem}{Problem}
\newtheorem{lemma}{Lemma}
\newtheorem{theorem}{Theorem}
\newcommand{\Bo}[1]{{\color{red} Bo: #1}}
% \newcommand{\QM}[1]{{\color{blue} QM: #1}}
\newcommand{\QM}[1]{{\color{blue}{#1}}}

\newcommand{\D}{\mathcal{T}}
\newcommand{\V}{\mathsf{V}}
\newcommand{\oR}{\mathcal{R}}
\newcommand{\MU}{\mathsf{U}}
\newcommand{\vats}{\mathsf{VFGS}}
\newcommand{\rand}{\mathsf{RAND}}
\newcommand{\full}{\mathsf{FULL}}
\newcommand{\avats}{\mathsf{VFGS}^{+}}
\newcommand{\cavats}{\mathsf{VFGS}^{+}\mathsf{CE}}
\newcommand{\sz}{\textsf{Shenzhen}}
\newcommand{\pt}{\textsf{Porto}}
\newcommand{\trim}{\vspace{-2mm}}

\newcommand{\stitle}[1]{\vspace*{0.4em}\noindent{\bf #1:\/}}
\newcommand{\sstitle}[1]{\vspace*{0.4em}\noindent{\bf #1\/.}}


\newcommand{\squishlist}{
	\begin{list}{$\bullet$}
		{ \setlength{\itemsep}{0pt}
			\setlength{\parsep}{3pt}
			\setlength{\topsep}{3pt}
			\setlength{\partopsep}{0pt}
			\setlength{\leftmargin}{1.2em}
			\setlength{\labelwidth}{1em}
			\setlength{\labelsep}{0.6em}
		}
	}
	\newcommand{\squishend}{
	\end{list}
}

% Defined by authors --------------------------------- Done

\begin{abstract}
Visualizing large-scale trajectory data is a core subroutine for many smart city applications, e.g., traffic management, urban planning, and route recommendation. The task is challenging due to high scalability requirement and severe visual clutter. Sampling can effectively mitigate both problems but may harm visual fidelity, i.e., generating visualizations that look different from the exact one. In this work, we propose sampling techniques that provide \textit{theoretically guaranteed visual fidelity} for large-scale trajectory data visualization. We first define a natural pixel-based \textit{fidelity loss function} to measure the difference between two visualizations. However, our analysis shows that it is NP-hard to sample a fixed-size subset of trajectories to minimize the loss function. Therefore, we then devise an approximate algorithm named \textit{$\vats$} with a suite of optimization techniques, which runs efficiently but still provides guaranteed fidelity loss. By taking data distribution and human perception characteristics into consideration, we further improve $\vats$ to an advanced algorithm named $\avats$ to tackle the visual clutter problem. Extensive experiments (i.e., case-, user-, and quantitative- studies) are also conducted on real-world trajectory datasets to verify the effectiveness and efficiency of our proposals.
%In addition, comprehensive user studies further illustrate the superiority of our proposals in various applications, e.g., traffic flow comparison, and reachable route inspection.
\end{abstract}

\begin{IEEEkeywords}
component, formatting, style, styling, insert
\end{IEEEkeywords}

% wari

\input{sections/Introduction.tex}

\input{sections/RelatedWork.tex}

\input{sections/ProblemFormulation.tex}

\section{Our Solution: $\vats$}\label{sec:sol}
In this part, we address the second technical challenge: \emph{how to devise an efficient sampling algorithm that provides guaranteed visual quality?}
Specifically, we first propose a visual quality-guaranteed sampling algorithm $\vats$ in Section~\ref{sec:greedy}. Next, we devise optimizations to improve the efficiency of $\vats$ in Section~\ref{sec:opt}.

\subsection{Visual quality-Guaranteed Sampling}\label{sec:greedy}
%Due to the hardness of Problem~\ref{prob:def}, the straight-forward solution is uniform random sampling $\rand$.
%This solution randomly selects $k$ trajectories from the dataset $\D$ and stores them in the result set $\oR$. The selected trajectories in $\oR$ are rendered as the visualization result.
%Obviously, uniform random sampling $\rand$ does not provide any guarantee on the visual quality of the sampled set.

Our visual quality-guaranteed sampling method $\vats$ is presented in Algorithm~\ref{alg:greedy}, which employs a greedy paradigm.
In particular, it finds the trajectory $tmp$ in $\D$ that maximizes $| tmp \cup \oR|$ at each iteration, as shown in Line~\ref{line:max} of Algorithm~\ref{alg:greedy}.
It terminates after $k=\alpha |\D|$ iterations and returns $\oR$ as the result set for rendering.

%\vspace{-2mm}
\begin{algorithm}
    \caption{$\vats(\D,k=\alpha |\D|)$} \label{alg:greedy}
    \begin{algorithmic}[1]
    \State Initialize result set $\oR \leftarrow \emptyset$
    \While{$|\oR| < k$}
        \State $tmp \leftarrow argmax_{t_i \in \D} |t_i \cup \oR|$ \label{line:max}
        \State $\oR \leftarrow \oR \cup \{ tmp \}$
    \EndWhile
    \State Return $\oR$
    \end{algorithmic}
\end{algorithm}


Interestingly, Algorithm~\ref{alg:greedy} provides provable visual quality guarantee for the returned result $\oR$, as stated in Theorem~\ref{the:ratio}.
%~\footnote{We omitted all the proofs of the theorems and lemmata in this work due to space reasons, and refer the interested readers to our technical report\cite{techreport}.}

\begin{theorem}~\label{the:ratio}
Define the visual quality of a sample set $\mathcal{S}$ as $f(\mathcal{S})=1-loss(\V(\D),\V(\mathcal{S}))$. Let the sized-k result produced by Algorithm~\ref{alg:greedy} be $\mathcal{R}$ and the sized-k solution of Equation~\eqref{eq:opp} be $\mathcal{R}^{\star}$, we have $f(\mathcal{R})\ge 0.632*f(\mathcal{R}^{\star})$.
\end{theorem}

Theorem~\ref{the:ratio} follows directly from the submodularity of the visual quality function $f(\mathcal{S})$ as it is well known that the greedy solution provides a $0.632$ approximation of the optimal solution for a submodular function. We show that submodularity of $f(\mathcal{S})$ as follows.      

\begin{lemma}[Submodularity]\label{lem:submodular}
	Define the contribution of a trajectory $t$ to the result set $\oR$ as $\Delta(\oR, t) = |V(\oR \cup t)| - |V(\oR)|$.
	Given a trajectory $t$ and two result sets $\oR,\oR^{'}$, if $\oR \subset \oR^{'}$, then $ \Delta(\oR, t) \geq \Delta(\oR^{'}, t)$.
\end{lemma}


%\REPORT{
\begin{proof}
	The contribution value of trajectory $t$ w.r.t. a given result set $\oR$ (i.e., $\Delta(\oR, t) = |V(\oR \cup t)| - |V(\oR)|$) is the newly covered pixels, which can be expressed as $|V(t)| - |V(\oR \cap t)|$.
	We have $t \cap \oR \subseteq  t \cap \oR^{'}$ as $\oR^{'}$ is a superset of $\oR$, which implies $|V(t)| - |V(t \cap \oR)| \geq |V(t)| - |V(t \cap \oR^{'})|$.
	Thus, it holds that $\Delta(\oR, t) = |V(\oR \cup t)| - |V(\oR)| \geq |V(\oR^{'} \cup t)| - |V(\oR^{'})|= \Delta(\oR^{'}, t)$.
\end{proof}
%}


%\begin{proof}
%The optimal solution of Problem~\ref{prob:def} covers $f(\mathcal{R}^{\star})$ pixels in $k$ iterations.
%Let $a_i$ be the number of newly covered pixels at the $i$-th iteration, $b_i$ is the total number of covered pixels up to the $i$-th iteration (i.e., $b_i = \sum_{j=1}^{i}a_i$),
%and $c_i$ be the uncovered pixels after $i$-th iteration (i.e., $c_i = f(\mathcal{R}^{\star})-b_i$).
%According to greedy paradigm, we can conclude the number of newly covered pixels at the $(i+1)$-th iteration is always greater than or equal to $\frac{1}{k}$ of the number of uncovered pixels after the $i$-th iteration, i.e., $a_{i+1} \geq \frac{c_i}{k}$.
%We prove Theorem~\ref{the:ratio} by proving $c_{i+1} \leq (1-1/k)^{i+1} \cdot f(\mathcal{R}^{\star})$.
%It holds $c_1 \leq (1-1/k) \cdot f(\mathcal{R}^{\star})$ as follows.
%\begin{align} \nonumber
%& a_1 \geq c_0 \cdot 1/k = 1/k \cdot f(\mathcal{R}^{\star}) \text{~~~as we concluded~~~} a_{i+1} \geq \frac{c_i}{k}\\ \nonumber
% \Leftrightarrow  & b_1 \geq 1/k \cdot f(\mathcal{R}^{\star})  \Leftrightarrow  -b_1 \leq - 1/k \cdot f(\mathcal{R}^{\star})  \text{~~~as~~~} a_1 = b_1\\ \nonumber
% \Leftrightarrow & f(\mathcal{R}^{\star}) - b_ 1 \leq f(\mathcal{R}^{\star}) - 1/k \cdot f(\mathcal{R}^{\star})  \Leftrightarrow  c_1 \leq (1-1/k) \cdot f(\mathcal{R}^{\star})
%\end{align}
%For inductive hypothesis assume $c_{i} \leq (1-1/k)^i \cdot f(\mathcal{R}^{\star})$. Thus,
%\begin{align} \nonumber
%& c_{i+1} = c_i - a_{i+1} \leq c_i - c_i/k = (1-1/k) \cdot c_i =  (1-1/k)^{i+1} \cdot f(\mathcal{R}^{\star})
%\end{align}
%
%Hence, it holds $c_k \leq (1-1/k)^{k} \cdot f(\mathcal{R}^{\star})$.
%It is equivalent to $b_k = f(\mathcal{R}) \geq (1 - (1-1/k)^{k}) \cdot f(\mathcal{R}^{\star}) \geq (1-1/e) \cdot f(\mathcal{R}^{\star}) \approx 0.632 \cdot f(\mathcal{R}^{\star})$.
%\end{proof}

Although Algorithm~\ref{alg:greedy} provides quality guarantee for the result set $\oR$, it has a high time complexity, which we show in the following analysis.  

 


%With the above theoretical analysis, Algorithm~\ref{alg:greedy} offers a visual quality-guaranteed sampling algorithm for the large-scale trajectory data visualization problem.
%However, as the time complexity analyzed in Lemma~\ref{lem:cost}, it is not scalable to large-scale trajectory datasets (e.g., millions of trajectories).


\begin{lemma}[Time Complexity]~\label{lem:cost}
	Given trajectory dataset $\D$ and an integer $k = \alpha |\D|$, the time complexity of Algorithm~\ref{alg:greedy} is $O(\alpha \cdot m \cdot |\D|^2)$, where $m$ is the maximum length of all trajectories in dataset $\D$.
\end{lemma}


\begin{proof}
	In each iteration,
	Algorithm~\ref{alg:greedy} computes the uncovered pixels of each trajectory in dataset $\D$ with $O(m)$ cost.
	As the dataset $\D$ has $O(|\D|)$ trajectories and Algorithm~\ref{alg:greedy} runs for $k = \alpha |\D|$ iterations, the total cost is $O(k \cdot m \cdot |\D|)=O(\alpha \cdot m \cdot |\D|^2)$.
\end{proof}


The high complexity of Algorithm~\ref{alg:greedy} hurts its scalability for large-scale trajectory datasets. For example, the \pt{} dataset contains 2.39 millions taxi trajectories, and the longest trajectory has 3,490 GPS points.
It takes 413.6 seconds for Algorithm~\ref{alg:greedy} to obtain a result set $\oR$ with sampling rate $0.1\%$.    Obviously, the running time is too long for interactive trajectory exploration.

\subsection{Heap-based Lazy Computation}\label{sec:opt}




\begin{figure}
	\centering
	\small
	\begin{tabular}{cc}
		\includegraphics[width=0.35\columnwidth]{pictures/1st}
		&
		\includegraphics[width=0.35\columnwidth]{pictures/2nd}
		\\
		(A) 1st iteration
		&
		(B) 2nd iteration
	\end{tabular}
	\caption{Lazy computing manner.} \label{fig:heap} %via the submodularity in Lemma~\ref{lem:submodular}
\end{figure}







Algorithm~\ref{alg:greedy} essentially adds the trajectory that maximizes $\Delta(\oR, t) = |V(\oR \cup t)| - |V(\oR)|$ to $\oR$ in each iteration. Lemma~\ref{lem:submodular} shows that the contribution of a trajectory (i.e., $\Delta(\oR, t)$) cannot increase when Algorithm~\ref{alg:greedy} runs for more iterations because $ \Delta(\oR^{'}, t) \le \Delta(\oR, t) $ for  $\oR \subset \oR^{'}$. Based on this property, we can use  $\Delta(\oR, t)$ calculated in the previous iterations to prune a trajectory $t$ from computation. If we have $\Delta(\oR, t)\le \Delta(\oR^{'}, t')$, in which $\oR$ and $\oR^{'}$ are a previous and the current sample set, respectively, we know that $t$ can not be added to the sample set in the current iteration.         

 

To implement this idea, we maintain a max-heap for the number of uncovered pixels of each trajectory and update the max-heap in a lazy fashion. 
That is, the contribution of a trajectory is computed only when necessary.
Figure~\ref{fig:heap}(a) shows a tiny max-heap example for the numbers of uncovered pixels of trajectories from $t_1$ to $t_7$ with result set $\oR=\emptyset$.
At the first iteration, the root node of the max-heap, $t_3$ in Figure~\ref{fig:heap}(A), is selected.
At the second iteration, the number of uncovered pixels of the root node $t_1$ is updated to 7 w.r.t. result set $\oR = \{t_3 \}$ (see the gray node in Figure~\ref{fig:heap}(B)).
Then $t_1$ is selected at the second iteration without computing the number of uncovered pixels for other trajectories, i.e., all white nodes in Figure~\ref{fig:heap}(B).
The reason is that the contributions of these trajectories are all less than 7 according to the submodularity in Lemma~\ref{lem:submodular}.

%In summary, the number of uncovered pixels in each trajectory will only be computed with the latest result set $\oR$ when it is necessary in the lazy computing manner,
%e.g., only $t_1$ will be updated at the 2nd iteration in Figure~\ref{fig:heap}.
%It reduces many unnecessary computations through the lazy updating manner, e.g., all white nodes did not update at the 2nd iteration in the above example.

%We then analyze the time complexity of Algorithm~\ref{alg:greedy} with lazy computing manner in Theorem~\ref{lem:lazy}.
%
%\begin{lemma}[Optimized Time Complexity]~\label{lem:lazy}
%Given trajectory dataset $\D$ and an integer $k= \alpha |\D|$, the time complexity of Algorithm~\ref{alg:greedy} with lazy computing manner is $O(\alpha \cdot m \cdot x |\D| \log |\D|)$, where $x$ is the number of contribution computations among all $k$ iterations and $x \ll |\D|$.
%\end{lemma}
%
%\begin{proof}
%It first takes $O(|\D|)$ time to construct the max-heap~\cite{cormen2009introduction}.
%It incurs $O( m \cdot x \log |\D|)$ cost to select the trajectory with maximum uncovered pixels at each iteration ($k$ iterations in total).
%Hence, the overall cost is $O(|\D| + k \cdot m \cdot t \log |\D|)$.
%\end{proof}


The efficiency of Algorithm~\ref{alg:greedy} is significantly improved with heap-based lazy computation.
To exemplify, Algorithm~\ref{alg:greedy} takes 413.6 seconds to return the results with sampling rate $0.1\%$ on the \pt{} taxi trajectory dataset. In contrast, our performance-optimized $\vats$ needs only 1.2 seconds.


\section{Advanced Approach: $\avats$}\label{sec:aa}

In the previous section, we presented the $\vats$ algorithm, which produces quality-guaranteed samples and runs efficiently. 
In this section, we focus on the third technical challenge: \emph{how to tackle the visual clutter problem in large trajectory visualization}?
In particular, we devise an advanced approach $\avats$ by considering
(i) trajectory data distribution, and (ii) human perception capability.
We elaborate (i) and (ii) by the examples in Figure~\ref{fig:delta}.


\stitle{Trajectory data distribution} Considering the \pt{} trajectory dataset, Figure~\ref{fig:delta}(A) is the visualization result of $\vats$ with sampling ratio $0.5\%$. It is obvious that the trajectories follow a non-uniform distribution, and there are some dense regions and sparse regions as illustrated by the rectangles in Figure~\ref{fig:delta}(A).

      
%Obviously, the real-world trajectory dataset is non-uniform distributed.
%For example, the trajectories in dense region are much more than those in the sparse region, as illustrated by the rectangles in Figure~\ref{fig:delta}(A).

\stitle{Human perception capability} Comparing Figures~\ref{fig:delta}(A) and (B), it is easier to tell their differences in the sparse regions than in the dense regions. This is because human beings have limited  perception capability, and hence two visualizations look indistinguishable if both of them contain a large number of points in the same area. This is exactly the case for the two dense regions in Figures~\ref{fig:delta}(A) and (B).       




\begin{figure}%[t]
	\centering
	\includegraphics[width=0.48\textwidth]{pictures/problemsolveing/delta_motivation.pdf}
	\caption{The advanced approach $\avats$ on \pt{} ($\alpha = 0.5\%$).} 
	\label{fig:delta}
\end{figure}

%(see Algorithm~\ref{alg:plus})

Based on the two observations above, we can improve  $\vats$ by delivering richer information in the sparse regions and reducing visual clutter in the dense regions. $\avats$ in Algorithm~\ref{alg:plus} achieves both objectives using a perception tolerance parameter $\delta$, which models the perception capability of humans at the highest level of details.
Specifically, if pixel $(x,y)$ in canvas is covered by the result set $\oR$ at the highest level,
the pixels around $(x,y)$, i.e., from $(x-\delta, y-\delta)$ to $(x+\delta, y+\delta)$, does not need to be covered as  they are in the perception tolerance of human beings. We can easily modify $\vats$ in Algorithm~\ref{alg:greedy} to incorporate the perception tolerance parameter $\delta$ as shown in Algorithm~\ref{alg:plus}. $\avats$ measures the contribution of each trajectory $t_i$ w.r.t the augmented set $\oR^{+}$ in Line~\ref{line:deltamax},
where $\oR^{+}$ includes both the selected trajectories and their tolerance pixels (in Line~\ref{line:delta}).      

  
%Taking the above two observations into consideration, we can further improve the returning result of visual quality-guaranteed sampling approach $\vats$ by
%delivering rich information at sparse regions and reducing visual clutter in dense regions.
%In this section, we devise the advanced approach $\avats$ (see Algorithm~\ref{alg:plus}) to achieve the above two objectives.
%In specific, we introduce perception tolerance parameter $\delta$ in $\avats$, which models the perception capability of humans at the highest level of details.
%In other words, suppose the pixel $(x,y)$ in canvas is covered by the result set $\oR$ at the highest level,
%the pixels around $(x,y)$, i.e., from $(x-\delta, y-\delta)$ to $(x+\delta, y+\delta)$, are not necessary to cover because they are in the perception tolerance of human beings.


%It measures the contribution of each trajectory $t_i$ w.r.t the selected trajectory set $\oR$'s augmented set $\oR^{+}$, i.e., the selected trajectories and their tolerance pixels.
%.
%The augmented set $\oR^{+}$ will be updated by the selected trajectory $tmp$ and its tolerance pixels set (in Line~\ref{line:delta}).

%
\begin{algorithm}
    \caption{$\avats(\D,k=\alpha |\D|,\delta)$} \label{alg:plus}
    \begin{algorithmic}[1]
    \State Initialize result set $\oR \leftarrow \emptyset$
    \State Initialize augmented result set $\oR^{+} \leftarrow \emptyset$
    \While{$|\oR| < k$}
        \State $tmp \leftarrow argmax_{t_i \in \D} | t_i  \cup \oR^{+} |$ \label{line:deltamax}
        \State $\oR \leftarrow \oR \cup \{ tmp \}$
        \State $\oR^{+} \leftarrow \oR^{+} \cup \mathsf{augment}(tmp, \delta)$\label{line:delta}
    \EndWhile
    \For{each $t$ in $\D$} \Comment{Representative encoding} \label{line:s}
        \State $tr \leftarrow argmin_{t_i \in \oR}{|t-\mathsf{augment}(t_i, \delta)|}$
        \State $tr.\mathsf{cnt}++$ \label{line:e}
    \EndFor
    \State Return $\oR$
    \end{algorithmic}
\end{algorithm}


%Interestingly, the visual clutter large trajectory visualization problem can be further reduced
%by encoding representative trajectories in $\oR$ (the returning result of $\avats$) with colors.
%In particular, $\avats$ selects the trajectory with the largest uncovered pixels by taking the perception tolerance capability of humans into account at each iteration,
%instead of only choosing the trajectory with the largest uncovered pixels in $\vats$ (Algorithm~\ref{alg:greedy}).


$\vats$ in Algorithm~\ref{alg:greedy} selects trajectories with good representativeness and some trajectories will not be included into the result set $\oR$ even though they have more uncovered pixels w.r.t. $\oR$. The reason is that their uncovered pixels are too close to the pixels in the selected trajectories (i.e., within the tolerance area of selected pixels). Take Figure~\ref{fig:zoom}(A) for example, suppose $\delta=1$ and trajectory $a$ was selected at the first iteration, the trajectory to select in the second iteration is $c$ instead of $b$ because almost all pixels in $b$ is in the tolerance area of $a$'s. $\avats$ also provides excellent visual quality at arbitrary zooming resolutions. This is because it considers the perception tolerance parameter $\delta$  at the highest zoom level. For example, the zoom level in Figure~\ref{fig:zoom}(A) is higher than that in Figure~\ref{fig:zoom}(B).
According to our elaboration, $\avats$ selects trajectory $a$ and $c$ for Figure~\ref{fig:zoom}(A).
When the area is zoomed out, as shown in Figure~\ref{fig:zoom}(b), trajectory $a$ and $c$ still captures the main sketch of the underlying dataset (as gray cells shown).
We will further elaborate this phenomenon by the experimental study in Section~\ref{sec:exp}.



The visual clutter problem for large-scale trajectory visualization can be further alleviated by encoding the representativeness of the trajectories in $\oR$ with colors. We define the representativeness of a trajectory $t_i$ in $\oR$ as the number of trajectories in the dataset $\D$ that has $t_i$ as its nearest neighbor in $\oR$. The distance between trajectory $t$ and $t_i$ is defined as the number of pixels in $t$ that can not be covered by the augmented pixels of $t_i$. We compute the representativeness of each trajectory in $\oR$ from Line~\ref{line:s} to Line~\ref{line:e} in Algorithm~\ref{alg:plus}. Figure~\ref{fig:delta}(C) shows the visualization result by encoding trajectories with larger representativeness with warmer colors. There are more details in the sparse regions compared with the $\vats$ result in Figure~\ref{fig:delta}(A), and we can identify the main roads in the dense region with very warm colors.

%Thus, the selected trajectories in the dense region are more representative than those in sparse region.

\begin{figure}[t]
	\centering
	\includegraphics[width=0.45\textwidth]{pictures/problemsolveing/one_to_many.pdf}
	\vspace{-2mm}
	\caption{$\avats$ at different zoom levels.}	\label{fig:zoom} %An illustration of
    \vspace{-2mm}
\end{figure}






%(1) Richer Information Delivering: details aware; so Arbitrary zooming resolutions
%(2) Popularity Embedding: visual clutter



%\subsection{One-to-many strategy}~\label{sec:one_to_many}
%Since we detect the covered pixels in the highest level, two trajectories may be very close to each other but share very few pixels, which will lead to more information loss in the low zoom view as figure~\ref{fig:one_to_many}.
%We next elaborate a ``one-to-many'' strategy to further optimize the visual quality of our proposed technique.
%Recalling we use the highest zoom level to define the pixel size in the canvas.
%Thus, our visual quality guaranteed sampling algorithm is zoom-level oblivious, e.g., it guarantees the visual quality of result set $\oR$ at every zoom level.
%However, users always do not use/need the highest zoom level in visualization applications.
%For example, Google map shows city and streets at zoom level 1 and 15, respectively~\footnote{\url{https://developers.google.com/maps/documentation/}}.
%Motivated by the above observation, we devise ``one-to-many'' strategy by introducing a visual tolerance parameter $\delta$ to optimize the visual quality for users.
%Specifically, ,
%the ``one-to-many'' strategy will ignore all the pixels around $(x,y)$ within $\delta$ offset distance, i.e., all pixels from $(x-\delta, y-\delta)$ to $(x+\delta, y+\delta)$ will be skipped.
%We will demonstrate the effectiveness of the visual tolerance $\delta$ in experimental evaluations.
%
%%https://developers.google.com/maps/documentation/maps-static/dev-guide#Zoomlevels
%\begin{figure}[t]
%	\centering
%	\includegraphics[width=0.4\textwidth]{pictures/problemsolveing/one_to_many.pdf}
%	\vspace{-5mm}
%	\caption{Resolution inconsistency}
%	\vspace{-5mm}
%	\label{fig:one_to_many}
%\end{figure}


%Specifically, $\avats$ incorporates a parameter $\delta$ during trajectory selection process in $\vats$ .
%In particular, we employ the parameter $\delta$ to model the end user's perception ability at the most high level of details.
%Surprisingly, our advance approach $\avats$ not only provides better visualization result when comparing with $\vats$ with the same sampling rate
%(e.g., Figure~\ref{fig:delta}(a) and (b) are the returning result of $\vats$ and $\avats$ respectively),
%but also embeds the popularity of selected trajectories by encoding the rest trajectories in the dataset in them,
%e.g., Figure~\ref{fig:delta}(c) is the visual result of $\avats$ with color encoded popularity.





\section{Experimental Evaluation}\label{sec:exp}

In this part, we first present several case studies of the results provided by $\avats$ in Section~\ref{sec:case} to demonstrate its good visualization quality. In Section~\ref{sec:user}, we conduct a comprehensive user study to test the effectiveness of $\avats$ on practical tasks including \textit{region center identification}, \textit{reachable route inspection} and \textit{traffic flow comparison}. In Section~\ref{sec:quality}, we quantitatively evaluate the visual quality and efficiency of $\avats$ and compare with well-designed baselines.     

\begin{table}
	\centering
	\small
	\caption{Statistics of the datasets used in the experiments}
	\begin{tabular}{|c|c|c|c|c|} \hline
		\textbf{Dataset} & \textbf{Trajectories} & \textbf{GPS Points} & \textbf{Maximum Length}\\ \hline
		\pt{}& 2,390,000 & 75,670,000 & 3,490 \\ \hline
		\sz{}& 3,070,000 & 53,530,000 & 2,268 \\ \hline
		\cd{}& 280,000 & 6,690,000 & 4,025 \\ \hline
	\end{tabular}	\label{tab:dataset}
\end{table}


\vspace{1mm}
\noindent\textbf{Experiment Settings.} We conduct the experiments 
using three real-world trajectory datasets: \pt{}, \sz{} and \cd{}.
\pt{}~\cite{pt} contains a total of 2.39 million taxi trajectories and 75.67 million of GPS points, and the longest trajectory has 3,490 GPS points.
\sz{}~\cite{sz} consists of 3.07 million taxi trajectories with 53.53 million GPS points, and the longest trajectory has 2,268 GPS points. 
\cd{}~\cite{cd} has \QM{0.28} million taxi trajectories and \QM{6.69} million GPS points, and the longest trajectory consists of \QM{4,025} GPS points. The statistics of the datasets are summarized in Table~\ref{tab:dataset}. All experiments are conducted on a machine with an Intel i7-8700 CPU, 24 GB memory and an NVIDIA GeForce GTX1080 GPU with 8 GB on-chip memory, running on Windows 10. All methods are implemented using Java 1.8, and the Processing 3 library~\cite{p3} is used for rendering. All timing results are measured in sing-thread mode. All datasets and source codes required to reproduce our results are available at~\cite{code}.

\vspace{1mm}

\noindent\textbf{Baselines.} We compare $\avats$ with three main baselines, i.e., $\mathsf{Full}$, $\mathsf{Random}$ and $\mathsf{DTW}$. $\mathsf{Full}$ visualizes all trajectories in the user selected region while $\mathsf{Random}$ selects trajectories in the user selected region at random for visualization. $\mathsf{DTW}$ is based on the DTW distance between trajectories~\cite{borcan2012improving} and designed by us to select trajectories with good diversity. Specifically, $\mathsf{DTW}$ samples the trajectory that maximizes the aggregate DTW distance to all remaining trajectories in each step. $\mathsf{DTW}$ takes days to run on our datasets because it needs to compute the DTW distance between all trajectory pairs and DTW distance computation has a quadratic complexity w.r.t. trajectory length.  In all comparisons, we ensure that $\mathsf{Random}$ and $\mathsf{DTW}$ use the same number of trajectories as $\avats$.                    




\subsection{Case Study}\label{sec:case}

We conduct case study on the \pt{} dataset to demonstrate the visualization quality of $\avats$ from the following aspects.

\begin{figure*}[t]
	\centering
	\includegraphics[width=0.95\textwidth]{pictures/case_study_icde/case_study_detail.pdf}
	\vspace{-1mm}
	\caption{Case study of the visualization quality of $\avats$ for two detail regions.}
	\label{fig:detailview}
	\vspace{-1mm}
\end{figure*}

\subsubsection{Overview visualization}

We analyze the visualization result for overview by considering the entire \pt{} dataset. 

\stitle{Consistently good visual quality at overview}
At zoom level 11, Figure~\ref{fig:overview}(A) is the visualization result of $\mathsf{Full}$ on the \pt{} dataset.
With a sampling rate $\alpha \!=\! 1\%$, Figures~\ref{fig:overview}(B), (C) and (E) are the visualizations produced by $\mathsf{Random}$, $\mathsf{DTW}$,   
and $\avats$, respectively. Comparing with Figure~\ref{fig:overview}(B) and (C), it is obvious that Figure~\ref{fig:overview}(E) is more similar to Figure~\ref{fig:overview}(A). In particular, Figure~\ref{fig:overview}(E) not only preserves the overall visual structure of the entire region but also keeps the details of cities that are far from the center (marked by the dashed cycles in the figure). However, the details of these cities are lost in Figure~\ref{fig:overview}(B) as $\mathsf{Random}$ is more likely to select trajectories in the dense region. $\mathsf{DTW}$ in Figure~\ref{fig:overview}(C) preserves more details than $\mathsf{Random}$ in the sparse regions but still cannot match the quality of $\avats$ in Figure~\ref{fig:overview}(E).

\stitle{Consistently good visual quality under different sampling rates}
Figures~\ref{fig:overview}(D) and (E) are the visualizations produced by $\avats{}$ with a sampling rate of $0.1\%$ and $1\%$. We can make two observations: (i) the larger the sampling rate, the better the visual quality, i.e., Figures~\ref{fig:overview}(E) is more similar to Figure~\ref{fig:overview}(A) compared with Figure~\ref{fig:overview} (D); (ii) the visualization of $\avats$ with a sampling rate of $0.1\%$ (i.e., Figure~\ref{fig:overview}(D)) looks even more appealing than the visualizations of $\mathsf{Random}$ and $\mathsf{DTW}$ with a sampling rate of $1\%$ (i.e., Figure~\ref{fig:overview}(B) and (C)) as Figure~\ref{fig:overview}(D) better captures the overall visual structure of Figure~\ref{fig:overview}(A).



\stitle{Color encoding effectively mitigates visual clutter} At zoom level 11 and with a sampling rate of $1\%$, Figures~\ref{fig:overview}(E) and (F) are the visualizations produced our $\avats$ and $\cavats$ (i.e., $\avats$ with color encoding), respectively.
Visual clutter is severe for $\mathsf{Full}$ (i.e., Figure~\ref{fig:overview}(A)) and $\avats$ (i.e., Figure~\ref{fig:overview}(E)) as many pixels are colored for the dense region in the center, which makes it difficult to identify the main routes. The visualization of $\cavats$ in Figure~\ref{fig:overview}(F) migrates this problem by encoding the trajectories with color, and it is easier to identify some prominent trajectories and busy routes. 

\vspace{1mm}
\subsubsection{Detail visualization}\label{sec:detail}




We analyze the visualization result for details by investigating two regions of interest in the \pt{} dataset in Figure~\ref{fig:detailview}. 
 
%We next present the effectiveness of our proposals with detail views by investigating two regions of interest in \pt{}, the dense region B and the sparse region C(shown as in Figure~\ref{fig:detailview}(8)).

%$\mathsf{Full}$
%$\mathsf{DTW}$
%$\mathsf{Random}$



\stitle{Dense region} At zoom level 15, region B in Figure~\ref{fig:detailview}(A) is the center of Porto and has the highest concentration of trajectories. Therefore, $\mathsf{Full}$ suffers from severe visual clutter and it is difficult to identify the road networks in in Figure~\ref{fig:detailview}(B1).  $\mathsf{Random}$ and $\mathsf{DTW}$ in Figure~\ref{fig:detailview}(B2) and (B3) reduce the visual clutter to some extent by sampling some trajectories. $\avats$ is more successful in reducing the visual clutter of $\mathsf{Full}$ and the main road networks are very clear in Figure~\ref{fig:detailview}(B4). In addition, $\avats$ preserves more micro structures of the trajectories than $\mathsf{Random}$ and $\mathsf{DTW}$, e.g., the circular route in the dashed circular region.                

% and causes serious visual clutter, as visualized in Figure~\ref{fig:detailview}(B1).
%For example, the circular structures of the main route(shown as the dashed circular region in Figure~\ref{fig:detailview}(B1)) is unclear.
%$\localavats$ alleviates the visual clutter by preserve the $1\%$ trajectories of the total regions but the clutter is still serious. Furthermore, $\avats$ performs better than $\localavats$ by preserving less trajectories and reduce the visual clutter. 


\stitle{Sparse region}
At zoom level 14, region C in Figure~\ref{fig:detailview}(A) contains the city of Casino Espinho and has fewer trajectories than the dense region in the center. In this case, the sampling methods need to keep the structures in the trajectories of $\mathsf{Full}$ in Figure~\ref{fig:detailview}(C1) to provide good visualization quality. However, $\mathsf{Random}$ and $\mathsf{DTW}$ in Figure~\ref{fig:detailview}(C1) and (C2) fail to meet this requirement, e.g., they show no trajectory in the dashed circle. In contrast, $\avats$ in Figure~\ref{fig:detailview}(C4) preserves these structures.          

%Region C includes the city of Casino Espinho at zoom level 14, which contains less trajectories than the center of Porto as the visualization result of full dataset shown in Figure~\ref{fig:detailview}(C1). 
%Given fix sampling rate $\alpha=1\%$, Figure~\ref{fig:detailview}(C2) indicates the visualization of $\localavats$. This visualization result misses a lot if detail information in this region, because the fix sampling rate preserves too few trajectories which is difficult to guarantee the visual quality.
%While $\avats$ in Figure~\ref{fig:detailview}(C3) performs much better than $\localavats$ as the sampling rate is automatically adjusted to according to the visual quality. In this visualization, the trajectory sketch of Casino Espinho is almost the same as it in Figure~\ref{fig:detailview}(C1), the visualized result of full dataset.

To sum up, the case study shows that $\avats$ effectively mitigates visual clutter with sampling and color encoding. Due to the quality-aware $\vats$ sampling algorithm, $\avats$ also provides better visualization quality than $\mathsf{Random}$ and $\mathsf{DTW}$ by preserving more micro structures in the trajectories.     

\subsection{User Study}\label{sec:user}

In this part, we conduct a comprehensive user study to evaluate the quality of visualizations generated by different methods.

\subsubsection{Settings}
We recruited ** participants with ** females, ** males, aged **-** with a mean of ** for the user study. The user study is conducted on the two larger datasets, i.e., \pt{} and \sz{}, and four methods are investigated, i.e., $\full$, $\rand$, $\mathsf{DTW}$ and $\avats$. We manually select 22 \textit{center points} in the two datasets and define 3 \textit{visualization scales} including:
large-scale region (zoom level less than 13), middle-scale region (zoom level between 13 and 15), small-scale region (zoom level more than 15). For each center point and visualization scale, we generate a \textit{comparable visualization group}, which includes one visualization generated by each of the 4 methods. This results in 66 comparable groups (22 center points $\times$ 3 scales) and 264 visualization results (66 comparable groups $\times$ 4 visualizations).

We are interested in the visual quality and visual clutter of the visualization results, and hence designed three tasks for a comparable group: \textit{T1}) rank the visualizations in the group from the highest visual quality to the lowest visual quality by 1-4;
\textit{T2}) rank the visualizations in the group from the least visual clutter to the most severe visual clutter by 1-4;
\textit{T3}) select the acceptable visualizations (multiple choices allowed) for analysis and choose the reason for those that are not selected, and we provide three reasons including ``severe visual clutter", ``poor visual quality" and ``others".  The user study system is a web-based platform, in which all visualizations are displayed with a resolution of 450*300.

\subsubsection{User study procedure}

When the participants enter the user study system, they are given a brief introduction and a tutorial on how to conduct the tasks to get familiar with the interface and tasks.
For each participant, we randomly select 16 comparable groups and generate 48 tasks.
For each comparable group, the 4 visualizations (\textit{without specifying generated by which method}) in one comparable group are shown on the same web-page and a participant is required to perform task T1, T2 and T3 by inspecting them.

%At last, the participants are interviewed to collect feedback after finishing the study and their answers are saved for result analysis.

\subsubsection{Result analysis}

\begin{figure*}[t]
	\centering
	\includegraphics[width=0.6\textwidth]{pictures/user_study/rank.png}
	%\vspace{-3mm}
	\caption{User study, rank distribution.}
	\label{fig:rank}
	%\vspace{-2mm}
\end{figure*}

The left plot of Figure~\ref{fig:rank} reports the quality ranking of 4 methods in T1. The results show that $\full$ ranks the 1st in most cases while $\avats$ usually ranks 1st or 2nd. In contrast, $\mathsf{DTW}$ and $\rand$ rank 3rd and 4th in most cases. This suggests that $\avats$ outperforms $\mathsf{DTW}$ and $\rand$ in visual quality. We also found that $\avats$ ranks 1st mainly for large-scale regions in which there are many trajectories.         

 
%Figure~\ref{fig:rank} left shows the ranking among different methods with x-axis indicating the ranking from the highest quality to the lowest quality and y-axis indicating the selecting number for the specific method and ranking. The visualization of $\full$ has the highest visual quality since it has no information loss according the quality definition. The selection of $\avats$ is mostly concentrated at the first and second, which is closely behind the $\full$. The selections of $\baseline$ and $\rand$ are contracted at the third and fourth respectively, which is confirmed both of these two methods performs worse than $\full$ and $\qtavats$ in guarantee the visual quality.

The right plot of Figure~\ref{fig:rank} reports the anti-visual clutter ranking in T2. The results show that $\full$ has the most severe visual clutter, ranking 4th in most cases. $\rand$ and $\mathsf{DTW}$ reduce visual clutter via sampling, and thus usually rank 2nd and 3rd. $\avats$ is the most successful in reducing visual clutter, ranking 1st in 155 out of the xxx cases.       

%Figure~\ref{fig:rank} right reports the ranking among different methods with x-axis indicating the ranking from the least clutter to the most clutter and y-axis indicating the total selecting number for the specific method and ranking. We observe that the $\qtavats$ is ranked at the first 155 times which is significantly more than the other methods. The number it ranks at the second, third and last are 68, 96 and 14. There is no doubt that the visualization of $full$ suffers the most sever visual clutter problem because 210 of 333 total answers rank $full$ at the fourth, while other methods can be used to help to reduce the visual clutter.

We report the frequency each method is selected as acceptable and why a method is not selected for T3 using bar chart in Figure~\ref{fig:accept_rate}. Each column corresponds to a method and from left to right, the lengths of the bars means ``favorable'', ``not favorable due to visual clutter'', ``not favorable due to poor visual quality'' and ``not favorable for other reasons''. The results show that $\avats$ is acceptable in 85\% of the cases, and the other methods have significantly lower acceptance rate than $\avats$.      




\begin{figure}[t]
	\centering
	\includegraphics[width=0.35\textwidth]{pictures/user_study/accept_rate.png}
	%\vspace{-5mm}
	\caption{User study, accept rate.}
	\label{fig:accept_rate}
	%\vspace{-8mm}
\end{figure}






\subsection{Quantitative Evaluation}\label{sec:quality}
In this part, we quantitatively evaluate our proposals on the three trajectory datasets in terms of both visual quality and efficiency.

\begin{figure*}[t]
	\centering
	\includegraphics[width=0.9\textwidth]{pictures/quantitative_study_icde/rate_quality.png}
	%\vspace{-3mm}
	\caption{Visual quality vs. sampling rates(T1).}
	\label{fig:sample_quality}
	%\vspace{-3mm}
\end{figure*}

\begin{figure*}[t]
	\centering
	\includegraphics[width=0.9\textwidth]{pictures/quantitative_study_icde/rate_rendertime.png}
	%\vspace{-3mm}
	\caption{Processing time vs. sampling rates.}
	\label{fig:rate_quality}
	%\vspace{-3mm}
\end{figure*}


\begin{figure*}[t]
	\centering
	\includegraphics[width=0.9\textwidth]{pictures/quantitative_study_icde/size_rendertime.png}
	%\vspace{-3mm}
	\caption{Response time vs. region size.}
	\label{fig:size_rendertime}
	%\vspace{-3mm}
\end{figure*}



\begin{figure*}[t]
	\centering
	\includegraphics[width=0.9\textwidth]{pictures/quantitative_study_icde/rate_algtime.png}
	%\vspace{-3mm}
	\caption{Algorithm running time vs. sampling rate.}
	\label{fig:rate_algtime}
	%\vspace{-3mm}
\end{figure*}
%\begin{figure}[t]
%	\centering
%	\includegraphics[width=0.5\textwidth]{pictures/quantitative_study_icde/sample_quality.png}
%	\vspace{-8mm}
%	\caption{Visual quality vs. sampling rates.}
%	\label{fig:sample_quality}
%	\vspace{-3mm}
%\end{figure}

%\begin{figure}[t]
%	\centering
%	\includegraphics[width=0.5\textwidth]{pictures/quantitative_study_icde/zoom_quality.png}
%	\vspace{-8mm}
%	\caption{Visual quality vs. zoom level.}
%	\label{fig:zoom_quality}
%	\vspace{-3mm}
%\end{figure}

\stitle{Visual quality} Figure~\ref{fig:sample_quality} reports the visualization quality (as defined in Equation~\eqref{eqref:loss}) of the methods under different sampling rate. We consider the entire region in each dataset for this experiment. The results show that both $\vats$ and $\vatss$ achieve significantly high quality than $\rand$ and $\mathsf{DTW}$ under the same sampling rate. This is because $\vats$ and $\vatss$ explicitly considers visual quality in the sampling process. Specifically, $\vatss$ has a quality close to 1 with a sampling rate less than 1\%. $\mathsf{DTW}$ has a higher quality than $\rand$ because it considers the diversity of trajectories. 

In Figure~\ref{fig:sample_quality}, we report the visualization time (i.e., the time to generate visualization using the sampled trajectories) for the methods under different sampling rates. We still consider the entire region in this experiment and the visualization time of $\full$ is included at the top of each figure for reference. The results show that all sampling methods achieve significantly shorter visualization time than $\full$, which confirms our observation that sampling is effective in improving visualization efficiency. Under the same sampling rate, $\vats$ and $\vatss$ take longer visualization time than $\rand$ and $\mathsf{DTW}$ because $\vats$ and $\vatss$ are more likely to select long trajectories for quality maximization. Combining Figure~\ref{fig:sample_quality} and~\ref{fig:sample_quality}, we can conclude that $\vatss$ can achieve high visualization quality with  short visualization time.   

We also eventuate the \textit{response time} of the $\avats$ framework under different quality guarantee and region size in Figure~\ref{fig:size_rendertime}. The response time is defined as the time taken to process the user query region to generate the visualization result. We contain the regions to be rectangles with a constant height/width ratio and measures the size of a region by dividing its height over the height of the entire region. Under each region size, we report the average response time of three typical regions, i.e., a dense region, a sparse region and a medium region. The results show that $\avats$ achieves a short response time (less than 1 second in all cases) for different region size and quality guarantees. The response time decreases rapidly when the region size shrinks as there are fewer trajectories in a smaller region. The response time required to achieve a high quality (e.g., 0.9) is not significantly longer than a low quality (e.g., 0.6) as quality improves quickly with the number of sampled trajectories. 

In Figure~\ref{fig:rate_algtime}, we report the running time of $\vatss$ with and without heap-based lazy computation under different sampling rate. The results show that the heap-based optimizations reduces the running time of $\vatss$ by 2-3 orders of magnitude. For the sampling rates we considered, $\vatss$ runs efficiently and can finish within 1 second for the entire dataset.                                  

 



%\stitle{Running time evaluation}

%\begin{figure}[t]
%	\centering
%	\includegraphics[width=0.5\textwidth]{pictures/quantitative_study_icde/sample_time.png}
%	\vspace{-8mm}
%	\caption{Time usage vs. sampling rates.}
%	\label{fig:sample_time}
%	\vspace{-3mm}
%\end{figure}

%\begin{figure}
%	\centering
%	\small
%	\begin{tabular}{cc}
%		\includegraphics[width=0.44\columnwidth]{pictures/tporto}
%		&
%		\includegraphics[width=0.44\columnwidth]{pictures/tshenzhen}
%		\\
%		(A) \pt{}
%		&
%		(B) \sz{}	
%	\end{tabular}
%	\vspace{-3mm}
%	\caption{Running time of $\vats$ w/ and w/o Lemma~\ref{lem:submodular}.}
%	\label{fig:cost}
%	\vspace{-6mm}
%\end{figure}
%\QM{unfininshed}
%We first conduct an experiment to evaluate the rendering cost by datasize. We vary the number of trajectories from 1000 to 1 million, which are randomly selected from \pt{} dataset. The experimental results are summarized in Table~\ref{tab:gpu}. We observe that the rendering cost is linear with the input data trajectories.
%
%We first report the running time of our $\vats$ algorithm in Figure~\ref{fig:cost} by varying the sampling rate from $0.01\%$ to $1\%$. The results show that $\vats$ is quite slow without the submodularity of contribution value, which agrees with Lemma~\ref{lem:submodular} in Section~\ref{sec:opt}.
%Then we shown the total time usage with sampling rate as Figure~\ref{fig:sample_time}. {*******}
%
%The optimized $\vats$ (e.g., $\vats$ with Lemma~\ref{lem:submodular}) outperforms $\vats$ by one to three orders of magnitudes on both datasets. The result show that running time of our $\vats$ algorithm is below 1 second in most cases. We have shown that $\vats$ provides good visualization performance with a low sampling rate (e.g., $0.1\%$ or $1\%$) in Section 6.1 and 6.2,  and Table~\ref{tab:gpu} suggests that the rendering latency scales almost linearly with dataset cardinality. By significantly reducing the dataset cardinality with sampling, $\vats$ can effectively reduces the rendering latency to make interactive visualization possible without sacrificing visual quality. For example, rendering the full $\pt{}$ dataset takes about \QM{34 seconds}, with a sampling rate of $1\%$, $\vats$ can bring down the rendering latency to less than 1 second.

\begin{table}
	\centering
	\small
	\caption{Visualization rendering cost}
	\begin{tabular}{|c|c|c|c|c|} \hline
		No. Trajs & GPS points & Mapping (s) & Rendering (s) & Total (s)\\ \hline
		1,000& 31,300 & 0.027 & 0.003 & 0.03 \\ \hline
		10,000& 31,6531 & 0.169 & 0.005 & 0.174\\ \hline
		100,000& 316,7120 & 1.701 & 0.057 & 1.758 \\ \hline
		1,000,000& 31,646,379 & 15.562 & 0.592 & 16.154 \\ \hline
	\end{tabular}	\label{tab:gpu}
\end{table}



\input{sections/ConclusionFutureWork.tex}


\bibliographystyle{IEEEtran}
\bibliography{ref}

\end{document}
